\section{Recursive completion trees}\label{sec:completion}

This section constructs the finite trees to which the entropy bound
will be applied.  Two small tensors supply the starting points;
a local degeneration preserves three readable color flags through
repeated completions, and a final deletion makes the entire tree
readable.

We use colors \(\colB,\colA,\colC\) owned by \(X,Z,Y\), respectively.
A \emph{three-flag tensor} is a W-partition, as in
Section~\ref{alg:four}, together with a conditional readable tree on
the support of each color.  Thus exactly one color is assigned to
each supported monomial, and its owning side reads whether that
color's flag is on.  The additional trees will let us continue
separating labels after a completion has identified a color word.

\subsection{Two starting tensors}

The second starting tensor is the five-term tensor \(T'\) from
\eqref{alg:five}, obtained from the \(q=1\) Coppersmith--Winograd
tensor \cite[Section~7]{cw}.  We retain its color flags and make the
conditional trees explicit.  The formulas below give short polynomial
witnesses for both starting budgets.

The first base has budget \(R=2\) and support
\[
 \colB=\{100\},\qquad \colA=\{001\},\qquad \colC=\{010\}.
\]
It is the coefficient of order one in
\[
 (x_0+t x_1)(y_0+t y_1)(z_0+t z_1)-x_0y_0z_0,
\]
whose constant coefficient is zero.
Its conditional trees are trivial.

The second base has budget \(R=3\) and support
\begin{equation}\label{eq:five-term-base}
 \colA=\{002,011\},\qquad
 \colB=\{200,110\},\qquad
 \colC=\{020\}.
\end{equation}
The triples give indices on \(X,Y,Z\), with unit coefficients.
The flag tests are \(z>0\), \(x>0\), and \(y=2\).
Within \(\colA\), the choice is read by \(Y,Z\); within \(\colB\),
it is read by \(X,Y\).

To verify the budget, put
\[
 P(t)=(x_0+t x_1+t^2x_2)
      (y_0+t y_1+t^2y_2)
      (z_0+t z_1+t^2z_2).
\]
The first coefficient of
\(P(t)+P(-t)-2P(0)\) is twice the sum of the six monomials with
index sum two; the family has rank at most three.
Now give weight one to \(x_1,z_1,y_0,y_2\) and weight zero to all
other variables.
The term \(101\) has total weight three and the other five terms
have weight one.
The leading part is exactly \eqref{eq:five-term-base}.

\subsection{Completion and termination}

\begin{definition}\label{def:completion}
A completion of size \(m\ge2\) with center \(i\) takes \(m\) copies
of a three-flag tensor and retains the color words in which at most
one distinct noncenter color occurs.
A pure-center word has new color \(i\); otherwise its new color is
the unique noncenter color \(j\) that occurs.
\end{definition}

\begin{lemma}\label{lem:completion}
Completion is a local leading-term degeneration.
It produces another three-flag tensor with conditional readable trees.
If the old tensor contains \(b\) base factors of budget \(R\), the new
one has \(mb\) such factors and budget \(R^{mb}\).
\end{lemma}

\begin{proof}
Use the local tests from Lemma~\ref{alg:completion}, with color \(i\)
as center.  On the center's owning side, give weight one to coordinates for which
all \(m\) old center flags are on.
On each noncenter's owning side, give weight one when any old flag of
that color is on.
All other weights are zero.
For a supported color word, the total weight is
\[
 \ind_{\{\text{all center}\}}+
 \sum_{j\ne i}\ind_{\{\text{some }j\}}.
\]
It equals one exactly on the desired words, and two if both noncenter
colors occur.
The same tests supply the new solo flags.

Conditional on new color \(j\ne i\), the word uses only \(i,j\)
and contains at least one \(j\).
The owning side of \(i\) reads its positions, and the owning side of
\(j\) reads the complementary positions.
Thus either side reads the complete color word.
After that word is known, apply the old conditional trees successively
in the slots.
For new color \(i\), the word is already determined and only the old
trees are needed.
All \(m\) copies are charged, including slots with deterministic color.
\end{proof}

To terminate, delete one color using its owning flag.
The choice between the two surviving colors is read by both their
owning sides.  Follow it with their conditional trees.
This is a readable tree to which \cref{thm:entropy} applies.

A \emph{script} consists of a base, a finite list of completions, and
a termination pair.
Its base-factor count is
\[
 b=\prod_{\text{completions}}m
\]
(the empty product is one), and its cost before the auxiliary tensors
of the readable tree are introduced is \(R^b\).
No conditional branch is exempt from this cost.
